\section{Sharp reservoir scales}
\label{sec:thresholds}

The main distinction is between three energy locations on the same
diverging scale and two narrow phase regions separated on a larger scale.
The following theorem needs only a weak energy limit.

\subsection{Three support points determine a quadratic threshold}

\begin{theorem}[Weak-support criterion]
\label{thm:three-support}
Suppose that, under the canonical reference law,
\begin{equation}
 \frac{E_N-a_N}{b_N}\ \Rightarrow\ X,
 \qquad b_N\to\infty,
 \label{eq:weak-energy-limit}
\end{equation}
and that the law of $X$ has at least three distinct points in its support.
For any prescribed exponents $c_N>0$,
\begin{equation}
 \bigl[\exists\hbox{ admissible }\mathcal E_N:
       \TV(Q_N,P_N)\to0\bigr]
 \quad\Longleftrightarrow\quad c_N\gg b_N^2.
 \label{eq:three-support-threshold}
\end{equation}
The centered choice~\eqref{eq:centered-calibration} suffices. The
necessity applies to every total-energy tuning, without a density,
moment, or local-limit assumption.
\end{theorem}

We first establish the geometric inequality used for necessity.

\begin{lemma}[Three-energy chord defect]
\label{lem:chord}
For $0<\theta<1$ and $q>0$, define
\begin{equation}
 f_\theta(q)=\log[(1-\theta)e^q+\theta]-(1-\theta)q.
 \label{eq:chord-function}
\end{equation}
Then
\begin{equation}
 f_\theta(q)\ge\frac{\theta(1-\theta)}{2e}
                  \min\{q^2,q\}.
 \label{eq:chord-lower-bound}
\end{equation}
If $g(e)=\beta_N e+c_N\log(\mathcal E_N-e)$ and
$e_1<e_2<e_3<\mathcal E_N$, set
\begin{equation}
 D=e_3-e_1,\qquad \theta=\frac{e_2-e_1}{D},\qquad
 q=\log\frac{\mathcal E_N-e_1}{\mathcal E_N-e_3}.
 \label{eq:three-point-parameters}
\end{equation}
The exact identities are
\begin{align}
 g(e_3)-g(e_1)&=\beta_ND-c_Nq,
 \label{eq:endpoint-identity}\\
 g(e_2)-(1-\theta)g(e_1)-\theta g(e_3)&=c_Nf_\theta(q).
 \label{eq:chord-identity}
\end{align}
\end{lemma}
\begin{proof}
We have $f_\theta(0)=f_\theta'(0)=0$ and
\[
 f_\theta''(q)=\frac{\theta(1-\theta)e^q}
 {[(1-\theta)e^q+\theta]^2}.
\]
For $0\le q\le1$, the denominator is at most $e^{2q}$, giving
$f_\theta''(q)\ge e^{-1}\theta(1-\theta)$. Integrating twice proves
\eqref{eq:chord-lower-bound} for $q\le1$. Convexity and
$f_\theta(0)=0$ imply that $f_\theta(q)/q$ is nondecreasing for $q>0$,
which proves it for $q\ge1$. For the identities, the linear energy
term cancels in the chord and
$\mathcal E_N-e_2=(1-\theta)(\mathcal E_N-e_1)
 +\theta(\mathcal E_N-e_3)$. Substitution gives both formulas.
\end{proof}

\begin{proof}[Proof of Theorem~\ref{thm:three-support}]
For sufficiency, use~\eqref{eq:centered-weight}. Weak convergence implies
$E_N-a_N=O_{P_N}(b_N)$. If $c_N\gg b_N^2$, the event
$|\beta_N(E_N-a_N)/c_N|\le1/2$ has probability tending to one. On it,
Taylor's theorem gives
\[
 \log W_N(E_N)
 =-\frac{\beta_N^2(E_N-a_N)^2}{2c_N}
 \left[1+O\!\left(\frac{|E_N-a_N|}{c_N}\right)\right]
 \longrightarrow0
\]
in probability. The uniform bound~\eqref{eq:global-centered-bound}
and Lemma~\ref{lem:normalization} prove convergence, and also show
admissibility for all large $N$. Changing finitely many choices of total
energy has no effect on the conclusion.

For necessity, let an arbitrary admissible tuning give convergence, and
write $R_N=L_N/Z_N$. Choose $\varepsilon_N\downarrow0$ slowly enough
that $P_N(|R_N-1|>\varepsilon_N)\to0$. On the complementary good set,
$R_N>0$ and $|\log R_N|\le2\varepsilon_N$ for all large $N$. Positivity
puts every good energy strictly below the reservoir cutoff.

Choose three bounded open intervals $I_1<I_2<I_3$, with positive gaps,
around three distinct support points of $X$. Each has positive limiting
lower probability by the Portmanteau theorem. Therefore each event
$\{(E_N-a_N)/b_N\in I_j\}$ intersects the good set for all large $N$.
Select one corresponding energy $e_{j,N}$ in each intersection. No
probability density or positive mass at an individual selected value is
required: the selection is from a nonempty event. The parameters in
\eqref{eq:three-point-parameters} obey
$D_N\asymp b_N$ and $\theta_N\in[\delta,1-\delta]$ for some fixed
$\delta>0$. The three log likelihoods are $o(1)$, so
\eqref{eq:endpoint-identity}--\eqref{eq:chord-identity} imply
\begin{equation}
 c_Nq_N=\beta_ND_N+o(1),\qquad
 c_Nf_{\theta_N}(q_N)=o(1).
 \label{eq:good-point-identities}
\end{equation}
The first quantity diverges. The chord lower bound then forces
$q_N\to0$: otherwise $c_N\min(q_N^2,q_N)$ is bounded below along a
subsequence by a positive multiple of $c_Nq_N\to\infty$.
It follows that $c_Nq_N^2\to0$. Finally,
\[
 \frac{D_N^2}{c_N}
 =\frac{c_Nq_N^2}{\beta_N^2}[1+o(1)]\longrightarrow0.
\]
Since $D_N\asymp b_N$, this is the claimed necessity.
\end{proof}

\begin{proposition}[The two-atom exception]
\label{prop:two-atoms}
If the canonical energy law has exactly two atoms $e_-<e_+$, with any
positive probabilities, then for every $c>0$ the choice
\eqref{eq:secant-calibration} reproduces the complete canonical state
law exactly.
\end{proposition}
\begin{proof}
The residual weight has equal positive values at both atoms, so its
normalized value is one $P$-almost surely. Apply
Lemma~\ref{lem:energy-tv}.
\end{proof}

For example, a nondegenerate Gaussian energy limit on the $\sqrt N$
scale gives $c_N\gg N$. Three or more positive-weight phase energies
in a limiting $E_N/N$ law give $c_N\gg N^2$. A continuous non-Gaussian
limit also satisfies the theorem, on whatever scale is established for
that model. A two-atom \emph{scaling limit} is different from an exactly
two-atom law: fluctuations hidden by the coarse scaling can impose a
larger requirement than Proposition~\ref{prop:two-atoms} reveals.

\subsection{Two phases with unequal energy scales}

The next result states precisely how fluctuations within a phase supply
the third relevant energy. Phase measures here and below are positive
probability measures. They need not have disjoint supports. They can also
be defined using labels on an enlarged probability space, provided $E_N$
is the actual subsystem energy throughout.

\begin{theorem}[Two-scale necessity]
\label{thm:two-scale-necessity}
Suppose the energy law admits the positive decomposition
\begin{equation}
 P_N^E=w_{-,N}P_{-,N}+w_{+,N}P_{+,N}+\delta_N R_N^{\mathrm{exc}},
 \quad w_{-,N}+w_{+,N}+\delta_N=1,
 \label{eq:positive-decomposition}
\end{equation}
where $\liminf_N w_{i,N}>0$ for $i=-,+$. Let
$\Delta_N=e_{+,N}-e_{-,N}\to\infty$ and
$s_N\to\infty$ with $s_N=o(\Delta_N)$. In at least one phase, say
the lower phase,
\begin{equation}
 (E_N-e_{-,N})/s_N\ \Rightarrow\ X_-
 \quad\hbox{under }P_{-,N},
 \label{eq:one-phase-weak-limit}
\end{equation}
where $\supp X_-$ has at least two distinct points. In the other phase,
assume $(E_N-e_{+,N})/\Delta_N\to0$ in probability. Then any admissible
total-energy tuning for which $\TV(Q_N,P_N)\to0$ must satisfy
\begin{equation}
 c_N\gg\Delta_Ns_N.
 \label{eq:two-scale-necessity}
\end{equation}
The same assertion holds with the phase roles interchanged.
\end{theorem}
\begin{proof}
The good set from the previous proof has probability tending to one
under each phase law: its complementary phase probability is bounded
by the unconditional complementary probability divided by $w_{i,N}$.
Choose two bounded separated neighborhoods of support points of $X_-$.
They supply good energies $e_{1,N}<e_{2,N}$ with
$e_{2,N}-e_{1,N}\asymp s_N$ and
$e_{j,N}=e_{-,N}+O(s_N)$. Concentration in the upper phase supplies
a good energy $e_{3,N}=e_{+,N}+o(\Delta_N)$; to make this selection
explicit, choose a deterministic sequence $r_N\downarrow0$ slowly
enough that $P_{+,N}(|E_N-e_{+,N}|\le r_N\Delta_N)\to1$.
Thus $D_N\asymp\Delta_N$ and
$\theta_N(1-\theta_N)\asymp s_N/\Delta_N$.

The identities~\eqref{eq:good-point-identities} still hold. Since
$\beta_N$ is bounded above and away from zero and $D_N\to\infty$,
they imply $q_N\asymp\Delta_N/c_N$. The chord bound is therefore at
least a positive constant times
\[
 c_N\frac{s_N}{\Delta_N}\min(q_N^2,q_N)
 \asymp\min\left\{\frac{\Delta_Ns_N}{c_N},s_N\right\}.
\]
It must tend to zero. Since $s_N\to\infty$, this forces
$\Delta_Ns_N/c_N\to0$. If the upper phase provides the two points,
the same argument has $1-\theta_N\asymp s_N/\Delta_N$.
\end{proof}

This obstruction is local to positive-probability phase neighborhoods.
It requires neither a Gaussian density nor control of the interphase
valley. In particular, extensive separation and nondegenerate weak
$\sqrt N$ fluctuations imply the necessary scale $c_N\gg N^{3/2}$.
The converse needs more information: under two-center calibration the
reservoir can amplify rare energies between the peaks by as much as
$\exp[O(\Delta_N^2/c_N)]$.

\subsection{A positive-phase sufficient criterion}

\begin{theorem}[Positive decomposition and exceptional mass]
\label{thm:positive-sufficiency}
Assume the positive decomposition~\eqref{eq:positive-decomposition},
with $\Delta_N\to\infty$, $s_N\to\infty$, and
$s_N=o(\Delta_N)$. Suppose that, for some fixed $\eta>0$, $M<\infty$,
and $N_0$,
\begin{equation}
 \sup_{N\ge N_0,\,i\in\{-,+\}}
 \E_{P_{i,N}}\exp\left(\eta\frac{|E_N-e_{i,N}|}{s_N}\right)\le M.
 \label{eq:phase-exponential-moment}
\end{equation}
There is a fixed constant $C$ depending only on bounds for $\beta_N$
such that the conditions
\begin{equation}
 c_N\gg\Delta_Ns_N,
 \qquad
 \delta_N\exp(C\Delta_N^2/c_N)\longrightarrow0
 \label{eq:mixture-sufficient-conditions}
\end{equation}
ensure $\TV(Q_N,P_N)\to0$ with the two-center calibration
\eqref{eq:secant-calibration}. Positive limiting phase weights and
local limit theorems are not required for this sufficient statement.
\end{theorem}
\begin{proof}
Let $W_N=e^{h_N}$, with value zero at and above the cutoff. Since
$c_N/\Delta_N\to\infty$, Lemma~\ref{lem:secant-bounds} applies.
The exponential moment implies tightness on the phase scale.
Together with~\eqref{eq:phase-flatness}, it gives $W_N\to1$ in each
phase probability. The tangent bound gives globally
\[
 W_N(e)\le
 \exp\left(\epsilon_N\frac{|e-e_{i,N}|}{s_N}\right),\qquad
 \epsilon_N=C\Delta_Ns_N/c_N\longrightarrow0.
\]
For large $N$, $2\epsilon_N\le\eta$, so
$\E_{P_{i,N}}W_N^2\le M$. This is uniform integrability, and therefore
$\E_{P_{i,N}}|W_N-1|\to0$ for each phase. On the exceptional component,
\[
 \delta_N\E_{R_N^{\mathrm{exc}}}|W_N-1|
 \le\delta_N\{1+\exp(C\Delta_N^2/c_N)\}\longrightarrow0.
\]
Adding the positive components and using
Lemma~\ref{lem:normalization} completes the proof.
\end{proof}

\begin{corollary}[A sharp two-phase criterion]
\label{cor:two-phase-iff}
Assume the positive decomposition, scale hypotheses, and moment bound
of Theorem~\ref{thm:positive-sufficiency}. Suppose additionally that
the phase weights have positive lower limits, at least one scaled
phase law has a weak limit with two distinct support points, and
\begin{equation}
 \delta_N\le\exp(-b\Delta_N/s_N)
 \quad\hbox{for some fixed }b>0.
 \label{eq:exceptional-surface-bound}
\end{equation}
Then a total-energy tuning with $\TV(Q_N,P_N)\to0$ exists if and only
if $c_N\gg\Delta_Ns_N$. In particular, if
$\Delta_N\asymp N$ and $s_N\asymp\sqrt N$, the sharp scale is
$c_N\gg N^{3/2}$.
\end{corollary}
\begin{proof}
The moment bound makes the other phase tight on scale $s_N$, hence
concentrated on the smaller-than-gap scale required by
Theorem~\ref{thm:two-scale-necessity}. For sufficiency,
$\Delta_N^2/c_N=o(\Delta_N/s_N)$ and $\Delta_N/s_N\to\infty$,
so~\eqref{eq:exceptional-surface-bound} implies the second condition
in~\eqref{eq:mixture-sufficient-conditions}.
\end{proof}

More generally, if $\delta_N\le e^{-t_N}$ with $t_N\to\infty$,
the sufficient condition
\begin{equation}
 c_N\gg\max\{\Delta_Ns_N,\Delta_N^2/t_N\}
 \label{eq:exceptional-general-rate}
\end{equation}
follows immediately. It is not a necessary condition expressed solely
in terms of the exceptional mass: an exceptional component near the
phase endpoints may receive much less amplification than the global
maximum used in the proof. Nor can a signed remainder in a partition
function expansion be substituted for the positive measure
$\delta_NR_N^{\mathrm{exc}}$. The exponential moment must concern the
energy actually exchanged with the reservoir.

\subsection{An alternative criterion allowing interfacial tails}
\label{subsec:density-tails}

The positive-phase exponential moment is one route to sufficiency.
A uniform bound on the full energy density gives another. We state the
bound explicitly, rather than deducing it from a local Gaussian limit.
Let $p_N$ be the density of $P_N^E$ with respect to Lebesgue measure and
let $\varphi_v$ denote the centered Gaussian density of variance $v$.
Assume
\begin{align}
 &\Delta_N/N\longrightarrow\ell>0,
 \label{eq:density-phase-gap}\\
 &\sqrt N\,p_N(e_{i,N}+\sqrt N z)
    \longrightarrow w_i\varphi_{v_i}(z)
    \quad\hbox{in }L^1([-R,R])\quad\hbox{for every }R<\infty,
 \label{eq:local-density-assumption}
\end{align}
where $w_i>0$, $w_-+w_+=1$, and $v_i>0$. On the interval between the
centers, set $x_N(e)=\min(e-e_{-,N},e_{+,N}-e)$. Suppose fixed
$A,a>0$ and $\alpha\in[1/2,1]$ give
\begin{equation}
 p_N(e)\le\frac{A}{\sqrt N}
 \exp\left[-a\min\left\{\frac{x_N(e)^2}{N},x_N(e)^\alpha\right\}\right],
 \qquad e\in[e_{-,N},e_{+,N}].
 \label{eq:interior-density-envelope}
\end{equation}
Constants carry the necessary fixed energy units. The minimum in this
bound permits a transition from a Gaussian cost to a subextensive
interfacial cost. No envelope is imposed outside the phase interval.

It is useful to give the sufficient argument for a general concave
residual weight; the exact physical reservoir is then a direct
application. Let $I_N$ be an interval of feasible energies, with $h_N$
concave and twice continuously differentiable in its interior, and set
$e^{h_N}=0$ outside $I_N$. Require that $I_N$ contain the two centers
and all their fixed $R\sqrt N$ neighborhoods for large $N$. Suppose
\begin{equation}
 h_N(e_{-,N})=h_N(e_{+,N})=0,\qquad
 0\le-h_N''(e)\le K\kappa_N
 \label{eq:concave-curvature-upper}
\end{equation}
between the centers and on each such neighborhood, with a constant
$K$ independent of $N$ and $R$.

\begin{proposition}[Sufficiency with an interior density envelope]
\label{prop:density-sufficiency}
Under~\eqref{eq:density-phase-gap}--\eqref{eq:concave-curvature-upper},
if $\kappa_NN^{3/2}\to0$, then
\begin{equation}
 \E_{P_N}|e^{h_N(E_N)}-1|\longrightarrow0.
 \label{eq:density-weight-convergence}
\end{equation}
The normalized reweighted law therefore converges in total variation.
\end{proposition}
\begin{proof}
Concavity and the endpoint values imply that the endpoint derivatives
are bounded in absolute value by $K\kappa_N\Delta_N$; for example,
integrate $h_N'$ between the zero endpoints and use its
$K\kappa_N$ Lipschitz bound. Hence, uniformly for $|z|\le R$,
\begin{equation}
 |h_N(e_{i,N}+\sqrt N z)|
 \le K\kappa_N\Delta_N\sqrt N R
       +\tfrac12K\kappa_NNR^2\longrightarrow0.
 \label{eq:general-local-flatness}
\end{equation}
Interpolation between the endpoints also gives, inside the interval,
\begin{equation}
 0\le h_N(e)\le\tfrac12K\kappa_N
 (e-e_{-,N})(e_{+,N}-e)
 \le C\kappa_NN x_N(e).
 \label{eq:general-interior-gain}
\end{equation}
Set $\lambda_N=\kappa_NN^{3/2}$. For
$R\sqrt N\le x\le\Delta_N/2$, the ratios of the last bound to the
two costs in~\eqref{eq:interior-density-envelope} satisfy
\[
 \frac{C\kappa_NNx}{x^2/N}\le C\lambda_N/R,
 \qquad
 \frac{C\kappa_NNx}{x^\alpha}
 \le C\lambda_N N^{1/2-\alpha}.
\]
For each fixed $R>0$, both tend to zero. The reservoir-weighted
interior mass outside the two phase windows is consequently bounded
for large $N$ by
\begin{equation}
 \frac{2A}{\sqrt N}\int_{R\sqrt N}^{\Delta_N/2}
 \exp\left[-\frac a2\min\left\{x^2/N,x^\alpha\right\}\right]\dd x.
 \label{eq:weighted-interior-tail}
\end{equation}
Since $e^{-a\min(u,v)/2}\le e^{-au/2}+e^{-av/2}$, the Gaussian part
has limiting upper bound $2A\int_R^\infty e^{-az^2/2}\dd z$, and the
other part tends to zero for fixed $R>0$ (its integral over the positive
half-line is finite, and there is an additional $N^{-1/2}$ factor).

Outside the phase interval, concavity gives $h_N\le0$ wherever feasible,
and the weight is zero elsewhere. These tails cannot be amplified.
Equation~\eqref{eq:local-density-assumption}, with the weights summing
to one, implies
\[
 \lim_{R\to\infty}\limsup_{N\to\infty}
 P_N\left(E_N\notin\bigcup_i
 [e_{i,N}-R\sqrt N,e_{i,N}+R\sqrt N]\right)=0.
\]
On each fixed phase window~\eqref{eq:general-local-flatness} makes the
weight uniformly approach one. On the remaining interior use
\eqref{eq:weighted-interior-tail} and the corresponding unweighted
tail; on the exterior use $|e^{h_N}-1|\le1$. Taking $N\to\infty$
first and then $R\to\infty$ proves~\eqref{eq:density-weight-convergence}.
\end{proof}

\begin{corollary}[Physical two-phase threshold with interfacial tails]
\label{cor:density-physical-iff}
Under~\eqref{eq:density-phase-gap}--\eqref{eq:interior-density-envelope},
an exact reservoir~\eqref{eq:bath-dos} can reproduce the canonical law
in total variation by some total-energy tuning if and only if
$c_N\gg N^{3/2}$.
\end{corollary}
\begin{proof}
For sufficiency choose~\eqref{eq:secant-calibration}. If
$c_N\gg N^{3/2}$, the residual reservoir energy is
$(c_N/\beta_N)[1+o(1)]$ uniformly between the centers and in each
fixed phase neighborhood. Thus~\eqref{eq:concave-curvature-upper}
holds with $\kappa_N=1/c_N$, and
Proposition~\ref{prop:density-sufficiency} applies.

For necessity, divide the energy line at the midpoint of the centers.
The local limits and their total weight one imply that these two
positive conditional phase measures have weights tending to $w_i$
and scaled energy laws converging weakly to $\Normal(0,v_i)$.
Indeed their mass in every fixed scaled compact interval has the stated
limit, and the limiting Gaussian mass tends to the entire phase weight
as the interval increases; no mass can escape on either side of the
midpoint. Theorem~\ref{thm:two-scale-necessity} with $s_N=\sqrt N$
then gives the required condition for any tuning.
\end{proof}

The density hypotheses are stronger than local coexistence limits.
They are assumptions to be verified for the specified microscopic
Hamiltonian and its energy, not consequences of identifying two peaks.
Proposition~\ref{prop:density-sufficiency} concerns continuous densities;
Theorems~\ref{thm:three-support}, \ref{thm:two-scale-necessity}, and
\ref{thm:positive-sufficiency} apply without change to discrete, continuous,
or mixed laws.
